% paper/sections/03c_levelset_mapping.tex
%
% §3 CLS 法の続き（03_levelset.tex からの分割）
% 内容：ψ-φ 関係・逆変換・曲率の数値計算と安定化
%
\subsection{\texorpdfstring{$\psi$}{ψ} と \texorpdfstring{$\phi$}{φ} の関係}
\label{sec:levelset_mapping}

\begin{tcolorbox}[defbox, title={$\psi$--$\phi$ 対応付けと解析的逆変換（ロジット関数）}]
Conservative Level Set 変数 $\psi$ と符号付き距離関数 $\phi$ は以下の関係にある~\cite{OlssonKreiss2005, Desjardins2008}：
\begin{equation}
  \psi(\bm{x},t) \equiv H_\varepsilon(\phi(\bm{x},t)), \qquad \bnabla\psi = \delta_\varepsilon(\phi)\,\bnabla\phi
\end{equation}
物性値補間 $\rho(\psi), \mu(\psi)$ には $\psi$ を直接用いる．
曲率 $\kappa$ は $\psi$ の微分値から直接計算できる
（定理~\ref{thm:curvature_invariance}，§~\ref{sec:curvature_direct_psi_cls}）．
分相 PPE（Pressure Poisson Equation；§~\ref{sec:split_ppe_motivation}）解法で使用する
HFE（Hermite 場延長法；§~\ref{sec:field_extension}）で
$\phi$ が必要な場合は逆変換（式~\eqref{eq:logit_inverse}）を用いる．

\medskip
\textbf{逆変換（$\psi \to \phi$）：}
CLS 法の移流計算は $\psi$ で行うが，幾何量計算には $\phi$ が必要である．
$H_\varepsilon(\phi) = 1/(1+e^{-\phi/\varepsilon})$ はシグモイド型関数であり，
その逆関数（ロジット関数）は\textbf{解析的に求まる}：
\begin{equation}
  \phi = H_\varepsilon^{-1}(\psi) = \varepsilon\ln\!\left(\frac{\psi}{1-\psi}\right)
  \label{eq:logit_inverse}
\end{equation}
$H_\varepsilon$ は単調増加関数（$H_\varepsilon' > 0$）であるから，解は一意に定まる．
唯一の実装上の注意は，$\psi \to 0$ または $\psi \to 1$ の\textbf{飽和領域}で
対数引数が $0$ または $\infty$ に近づき浮動小数点オーバーフローが生じる点である．
この場合は $|\phi| = \phi_{\max}$ に飽和させる（§~\ref{sec:newton_inversion} 参照）．
\end{tcolorbox}

\subsubsection{逆変換：解析的逆関数の性質}
\label{sec:newton_inversion}

式~\eqref{eq:logit_inverse} の逆変換は\textbf{ロジット関数}
$\phi = \varepsilon\ln(\psi/(1-\psi))$
によって点ごとに $\Ord{1}$ で実行できる．
$H_\varepsilon$ は狭義単調増加であるため逆関数は一意に定まり，
残差 $H_\varepsilon(\phi) - \psi$ は代数的にゼロである（ニュートン反復は不要）．

飽和領域（$\psi \to 0$ または $\psi \to 1$）では対数引数の特異性により
数値精度が低下するが，この領域は界面から遠く（$|\phi| \gg \varepsilon$），
曲率計算に寄与しない．
本稿ではロジット逆変換を\textbf{主手法}とし，
飽和クランプ $|\phi| \le \phi_{\max}$ でオーバーフローを防ぐ
（実装手順は本稿の 7-step フロー Step 3 に位置する）．

\paragraph{補助用途としての Newton 反復．}
診断や収束性の検証など補助目的では，以下の Newton 更新式が利用できる：
\begin{equation}
  \phi^{(k+1)} = \phi^{(k)} - \frac{H_\varepsilon(\phi^{(k)}) - \psi_i}{\dfrac{1}{\varepsilon}H_\varepsilon(\phi^{(k)})\bigl(1-H_\varepsilon(\phi^{(k)})\bigr)}
  \label{eq:newton_inversion}
\end{equation}
この反復の二次収束は，根に十分近い初期値（目安 $|\phi^{(0)}|\lesssim \varepsilon$）を用い，
$H_\varepsilon'(\phi)$ が小さすぎない局所領域でのみ成り立つ（付録~\ref{app:newton_quad_conv}）．
飽和クランプ単独では大域的な二次収束を保証しないため，
主計算経路としてはロジット逆変換が優先される．

\begin{tcolorbox}[mybox, title={ロジット逆変換 vs. Sussman 再初期化}]
\begin{center}
\renewcommand{\arraystretch}{1.3}
\begin{tabular}{lll}
\hline
 & \textbf{ロジット逆変換（本稿）} & \textbf{Sussman 再初期化} \\
\hline
計算式 & $\phi = \varepsilon\ln(\psi/(1-\psi))$（解析的） & $\partial\phi/\partial\tau + \mathrm{sgn}(\phi_0)(|\bnabla\phi|-1)=0$ \\
コスト & $\Ord{N}$（対数1回） & $\Ord{N \cdot N_\tau}$（PDE 反復求解） \\
精度 & $H_\varepsilon^{-1}(\psi)$ の厳密な逆変換 & $|\bnabla\phi|=1$ を修正するが $\psi$ との整合性なし \\
\hline
\end{tabular}
\end{center}
CLS 法では $\psi$ が基本変数であり，ロジット逆変換は $\psi$--$\phi$ の整合性を代数的に保つ．
\end{tcolorbox}

% ═══════════════════════════════════════════════════════════════════════════════
\subsection{\texorpdfstring{$\psi$}{ψ} 微分による曲率の直接計算：ロジット逆変換の省略}
\label{sec:curvature_direct_psi_cls}
% ═══════════════════════════════════════════════════════════════════════════════

第~\ref{sec:governing}章の定理~\ref{thm:curvature_invariance} により，
曲率は $\phi$ と $\psi$ のいずれから計算しても同一であることが示された．
本節では，この不変性を CLS 法に適用し，
ロジット逆変換（§~\ref{sec:newton_inversion}）を\textbf{省略}して
$\psi$ の微分値から曲率を直接計算する具体的手法を述べる．

\begin{tcolorbox}[resultbox, title={CLS 法における $\psi$ 微分による曲率の直接計算}]
式~\eqref{eq:curvature_2d} において $\phi$ を $\psi$ で置き換えた
\begin{equation}
  \kappa = -\frac{
    \psi_y^2\,\psi_{xx}
    - 2\psi_x\psi_y\,\psi_{xy}
    + \psi_x^2\,\psi_{yy}
  }{\bigl(\psi_x^2 + \psi_y^2\bigr)^{3/2}}
  \label{eq:curvature_psi_2d}
\end{equation}
は連続・非飽和領域では式~\eqref{eq:curvature_2d} と同一の曲率を与える
（定理~\ref{thm:curvature_invariance}）．
CCD 法は $\psi$ の解法過程で $\psi$, $\psi'$, $\psi''$ を
同時に $\Ord{h^6}$ 精度で算出するため，
式~\eqref{eq:curvature_psi_2d} の全ての入力が追加計算なしに利用可能である．

\medskip
\noindent\textbf{従来法（$\psi \to \phi \to \kappa$）に対する利点：}
\begin{enumerate}
  \item ロジット逆変換（式~\eqref{eq:logit_inverse}）の\textbf{完全な省略}：
        飽和領域（$\psi \to 0, 1$）での対数関数の特異性による誤差増幅が
        原理的に排除される．
  \item 逆変換後の $\phi$ に対する差分演算が不要となり，
        $\psi$ の滑らかな関数構造がそのまま保存される．
  \item 再初期化による Eikonal 条件 $|\bnabla\phi| \approx 1$ の維持が
        曲率精度の前提条件ではなくなる：
        ただし，$\psi$ が界面近傍で滑らかな単調変換として保たれていることが前提であり，
        プロファイルが強く歪んだ場合は再初期化や品質監視が必要である．
\end{enumerate}
\end{tcolorbox}

% ─────────────────────────────────────────────────────────────────────────────
\subsubsection{\texorpdfstring{数値的安定性：$|\bnabla\psi|$ の挙動と適用領域}{数値的安定性：|∇ψ| の挙動と適用領域}}
% ─────────────────────────────────────────────────────────────────────────────

式~\eqref{eq:curvature_psi_2d} の分母
$(\psi_x^2+\psi_y^2)^{3/2}$ は，$\psi$ のプロファイル形状に依存する．
界面近傍（$\psi \approx 0.5$）ではシグモイド勾配が最大値
$|\bnabla\psi|_{\max} = 1/(4\varepsilon)$（$|\bnabla\phi|=1$ のとき）
をとり，安定な除算が可能である．
一方，$\psi \to 0$ または $\psi \to 1$ の領域では
$|\bnabla\psi| \to 0$ となるが，
これらの領域は界面から遠く，表面張力項の評価に曲率が寄与しない．
したがって，実装上は以下のハイブリッド戦略をとる：
\begin{itemize}
  \item \textbf{界面近傍}（$\psi_{\min} < \psi < 1-\psi_{\min}$，
        推奨 $\psi_{\min}=0.01$）：
        式~\eqref{eq:curvature_psi_2d} を適用し，高精度な曲率を直接計算する．
  \item \textbf{界面遠方}（$\psi \le \psi_{\min}$ または
        $\psi \ge 1-\psi_{\min}$）：
        $\kappa = 0$ とする（表面張力寄与なし）．
\end{itemize}
この閾値選択に対するロバスト性は高い：
表面張力項 $\kappa\,\bnabla\psi/We$ において
$|\bnabla\psi| \to 0$ の領域では $\kappa$ が不定となっても
積 $\kappa\,|\bnabla\psi|$ は有界に留まるためである．

% ─────────────────────────────────────────────────────────────────────────────
\subsubsection{実効界面厚み \texorpdfstring{$\varepsilon_{\textup{eff}}$}{ε\_eff} によるプロファイル品質監視}
% ─────────────────────────────────────────────────────────────────────────────

$\psi$ による直接曲率計算では $\varepsilon$ の値は不要であるが，
数値拡散や再初期化の不完全さによるプロファイル歪みを
\textbf{診断}する目的で，各格子点における実効界面厚みを定義できる：
\begin{equation}
  \varepsilon_{\textup{eff},i}
  = \frac{\psi_i(1-\psi_i)}{|\bnabla\psi|_i}
\end{equation}
この定義は，シグモイド微分公式
$|\bnabla\psi| = \frac{1}{\varepsilon}\psi(1-\psi)\,|\bnabla\phi|$
と Eikonal 条件 $|\bnabla\phi| \approx 1$ から導かれる．
界面近傍で $\varepsilon_{\textup{eff}} \approx \varepsilon$（設計値）であれば
プロファイルが正常に維持されていることを意味し，
$\varepsilon_{\textup{eff}} \gg \varepsilon$ は
数値拡散による過度な界面幅増大を示す診断指標となる．

% ═══════════════════════════════════════════════════════════════════════════════
\subsection{曲率の数値的安定性}
\label{sec:curvature_stab}
% ═══════════════════════════════════════════════════════════════════════════════

$\psi$ 直接法（式~\eqref{eq:curvature_psi_2d}，定理~\ref{thm:curvature_invariance}）を用いる場合，
曲率計算は Eikonal 条件 $|\bnabla\phi|=1$ に依存しない．
適用域は系~\ref{cor:cls_curvature_domain} に従い，
界面近傍（$\psi_{\min} < \psi < 1-\psi_{\min}$，推奨 $\psi_{\min}=0.01$）で
式~\eqref{eq:curvature_psi_2d} を適用し，
界面遠方（飽和領域，$|\bnabla\psi| \to 0$）では $\kappa = 0$ とする．
CCD 法（第~\ref{sec:CCD}章）は解法過程で $\psi, \psi', \psi''$ を
同時に $\Ord{h^6}$ 精度で算出するため，追加計算は不要である．

$\phi$ 経由で曲率を計算する場合は，再初期化の不完全性により
分母 $|\bnabla\phi|$ が局所的に小さくなるリスクがある．
この場合は分母クランプ（式~\eqref{eq:curvature_clamped}，$\varepsilon_{\text{norm}} = 10^{-3}$）で
$\kappa$ を有界に保つ．曲率精度の第一の保証は正則化ではなく
再初期化の品質（Eikonal 条件の維持）にある——
Eikonal 誤差 $\delta \ll 1$ のとき曲率への影響は $\Ord{\delta^2}$ に留まる．

\medskip
以上で CLS 法の基本定式化（保存形移流・再初期化・ψ–φ 写像・曲率の直接計算）が整った．
次節（§~\ref{sec:ridge_eikonal}）ではトポロジー変化の連続化を扱う
Ridge–Eikonal ハイブリッド再初期化を導入し，第~\ref{sec:levelset}章を完結させる．
Compact Combined Difference（CCD）法の設計原理と離散化は
続く第~\ref{sec:CCD}章で述べる．
